\documentclass[12pt]{article}
\usepackage[margin=1in]{geometry}
\usepackage{setspace}
\usepackage{parskip}
\usepackage{enumitem}
\usepackage{titlesec}
\usepackage{xcolor}
\usepackage{mdframed}
\usepackage{hyperref}
\usepackage{amsmath}

% -------------------------------------------------------
% Styling
% -------------------------------------------------------
% Referee 1 � blue
\definecolor{r1color}{RGB}{30, 80, 150}
\definecolor{r1bg}{RGB}{240, 245, 255}
% Referee 2 � teal
\definecolor{r2color}{RGB}{0, 110, 100}
\definecolor{r2bg}{RGB}{235, 248, 245}

\newmdenv[
  backgroundcolor=r1bg,
  linecolor=r1color,
  linewidth=1.5pt,
  leftline=true, rightline=false, topline=false, bottomline=false,
  innerleftmargin=10pt, innerrightmargin=10pt,
  innertopmargin=6pt, innerbottommargin=6pt,
  skipabove=8pt, skipbelow=4pt
]{r1quote}

\newmdenv[
  backgroundcolor=r2bg,
  linecolor=r2color,
  linewidth=1.5pt,
  leftline=true, rightline=false, topline=false, bottomline=false,
  innerleftmargin=10pt, innerrightmargin=10pt,
  innertopmargin=6pt, innerbottommargin=6pt,
  skipabove=8pt, skipbelow=4pt
]{r2quote}

% Section formatting � colour changes per referee (set before each \section)
\titleformat{\section}{\large\bfseries\color{currcolor}}{}{0em}{}[\titlerule]
\titleformat{\subsection}{\normalsize\bfseries}{}{0em}{}

% Convenience commands
\newcommand{\response}{\par\smallskip\noindent\textbf{Response:}\ }

% -------------------------------------------------------
% Document metadata
% -------------------------------------------------------
\newcommand{\manuscriptid}{YGAME-D-25-00181}
\newcommand{\papertitle}{Preferences for Warning Signal Quality: Experimental Evidence}
\newcommand{\journalname}{[Journal Name]}
\newcommand{\submissiondate}{\today}

% -------------------------------------------------------
\begin{document}

\begin{center}
  {\LARGE\bfseries Response to Referee Reports}\\[6pt]
  {\large Manuscript \manuscriptid}\\[4pt]
  {\normalsize \textit{\papertitle}}\\[4pt]
  {\normalsize \journalname\ $\cdot$ \submissiondate}
\end{center}

\vspace{1em}
\hrule
\vspace{1em}

We thank both referees for their thorough and constructive reviews. The revised
manuscript is stronger for it. Below we address each point in turn. Referee comments
are shown in shaded boxes ({\color{r1color}blue} for Referee~1,
{\color{r2color}teal} for Referee~2); our responses follow in plain text.

\tableofcontents
\newpage

% ===============================================================
\colorlet{currcolor}{black}
\section{Associate Editor}
% ===============================================================

The Associate Editor has no additional comments.

% ===============================================================
\newpage
\colorlet{currcolor}{r1color}
\section{Referee 1}
% ===============================================================

We thank Referee~1 for the careful reading of both the original submission and the
revision.

%----------------------------------------------------------------
\subsection{R1.1\quad Incorrect and Inconsistent Definitions of FP and FN Costs}

\begin{r1quote}
The paper now emphasizes the distinction between false-positive/false-negative rates and
false-positive/false-negative costs. This is a useful clarification. However, the definitions
given in the model and in Section~5.2 are not correct. The paper defines FP and FN costs
on page~6 as $\pi(1-P_{01})c$ and $\pi P_{01}L$, respectively, and states that WTP is
equally affected by both. There are two problems.

\medskip
\textit{FP cost:} The expression $\pi(1-P_{01})c$ is the expected cost of true positive
protection, not the cost of false positives. The FP cost should be $(1-\pi)P_{10}c$,
consistent with equation~(4) where $\partial b^*/\partial P_{10} = -(1-\pi)c$.

\medskip
\textit{FN cost:} The paper defines it as $\pi P_{01}L$ (gross expected loss). An
alternative, $\pi P_{01}(L-c)$, captures the incremental cost: under a true-positive
signal the agent would have paid $c$, so the net cost of a miss is $L-c$, consistent
with equation~(5) where $\partial b^*/\partial P_{01} = -\pi(L-c)$.

\medskip
This matters for Table~5. The paper should explicitly define the regression variables,
map them to equation~(3), and verify that the definitions do not affect the reported
results.
\end{r1quote}

\response [Your response here.]

%----------------------------------------------------------------
\subsection{R1.2\quad Interpretation of Table~5 and Figure~5}

\begin{r1quote}
The abstract states that subjects ``disproportionately reduce their demand for signals
with high false negative rates for rare events.'' This is imprecise.

Under the risk-neutral null, all Table~5 coefficients should be zero; positive coefficients
mean WTP does not fall \emph{enough} as error costs rise. The actual finding is that
subjects selectively under-react to the \emph{dominant} cost type:
\begin{itemize}
  \item Low priors $\Rightarrow$ FP costs dominate $\Rightarrow$ subjects fail to adjust for FP costs.
  \item High priors $\Rightarrow$ FN costs dominate $\Rightarrow$ subjects fail to adjust for FN costs.
\end{itemize}
Figure~5 corroborates this: empirical sensitivity coefficients are flat across priors
whereas theory predicts sharp divergence. The finding is not absolute over-penalisation
of FN for rare events but a failure to \emph{differentially} weight FP and FN as the
prior shifts. The abstract and conclusion should be rewritten accordingly, and the paper
should use ``false alarm'' and ``missed detection'' alongside FP/FN notation.
\end{r1quote}

\response [Your response here.]

%----------------------------------------------------------------
\subsection{R1.3\quad Mechanism Claim Should Be Softened}

\begin{r1quote}
Section~5.3 is much improved, but I remain unconvinced that the data uniquely identify a
failure to distinguish FP and FN errors. The same pattern could reflect a total-error-rate
heuristic, a ``dishonest gremlin'' heuristic, probability weighting, noisy Bayesian
reasoning, task complexity, or some combination.

A more defensible conclusion: many subjects use a simplified total-error or
total-dishonesty heuristic rather than separately processing how FP and FN rates interact
with the prior. Relatedly, Section~4 should not be framed only as reassurance that
subjects understood the task � the BP, IP, and BE patterns are part of the behavioral
phenomenon itself.
\end{r1quote}

\response [Your response here.]

%----------------------------------------------------------------
\subsection{R1.4\quad Proposition~2 Wording}

\begin{r1quote}
Ensure the wording explicitly states that the proposition concerns the ratio of
sensitivities $(\partial b/\partial P_{01})/(\partial b/\partial P_{10})$, not absolute
levels. A brief economic interpretation would help: risk-averse agents are relatively
more intolerant of missed detections but more forgiving of false alarms compared with
risk-neutral agents.
\end{r1quote}

\response [Your response here.]

%----------------------------------------------------------------
\subsection{R1.5\quad Appendix~B: Proof Numbering}

\begin{r1quote}
The main text has two propositions but the appendix has three proofs (B.1--B.3).
Section~B.2 proves a result for imprudent agents that no longer corresponds to any
proposition in the main text. B.2 should be dropped or the main text updated to match.
\end{r1quote}

\response [Your response here.]

%----------------------------------------------------------------
\subsection{R1.6\quad Appendix~B.3: Algebra Error}

\begin{r1quote}
In equations~(13) and~(14) the denominator of the last expression is $(1-\pi)L$; it
should be $(1-\pi)c$:
\[
  \frac{\pi}{1-\pi}\cdot\frac{u'(\xi_1)(L-c)}{u'(\xi_2)\,c}
  = \frac{\pi(L-c)}{(1-\pi)\,c}\cdot\frac{u'(\xi_1)}{u'(\xi_2)}.
\]
\end{r1quote}

\response We thank the referee for catching this typo. The denominator in
equations~(13) and~(14) has been corrected from $(1-\pi)L$ to $(1-\pi)c$. The stated
result is unaffected.

%----------------------------------------------------------------
\subsection{R1.7\quad Cross-Reference Error}

\begin{r1quote}
Page~26 references ``Appendix Table~A10,'' but the table with matching content (WTP
deviations with IP bias controls) is numbered A8.
\end{r1quote}

\response [Your response here.]

%----------------------------------------------------------------
\subsection{R1.8\quad Typos}

\begin{r1quote}
\begin{itemize}[noitemsep]
  \item Page~7: ``Equations 6-s 7'' $\to$ ``Equations 6--7.''
  \item Page~15, fn.~10: missing comma between $\hat{\pi}=0.43$ and $\hat{\theta}=0.25$.
  \item Page~18: ``the value for for a prior of 0.2'' � duplicated word.
  \item Page~27: ``the observed sensitivity patterns emerges'' $\to$ ``emerge.''
  \item Page~27: ``characteristcs'' $\to$ ``characteristics.''
\end{itemize}
\end{r1quote}

\response All five typos have been corrected. We thank the referee for the careful
reading.

% ===============================================================
\newpage
\colorlet{currcolor}{r2color}
\section{Referee 2}
% ===============================================================

We thank Referee~2 for the detailed and constructive comments on the revision.

%----------------------------------------------------------------
\subsection{R2.1\quad Treatments and Treatment Assignment}

\begin{r2quote}
\textbf{(a)} The ``Treatments and Treatment Assignment'' subsection states that each
subject is randomly allocated to a pair of priors but does not explain how information
structures are assigned.

\textbf{(b)} The text notes that ``the first treatment was always the one with a fully
truthful information structure,'' suggesting information structures are \emph{not}
randomly assigned. The authors should provide a general description and justification.

\textbf{(c)} Because information structures do not appear to be randomly assigned,
sample sizes across treatment cells (prior $\times$ information structure) likely
differ. The authors should report sample sizes by treatment cell, e.g.\ in an appendix
table ($4\times12=48$ cells).
\end{r2quote}

\response [Your response here.]

%----------------------------------------------------------------
\subsection{R2.2\quad Description and Groupings in Tables~2, 3, and A4}

\begin{r2quote}
\textbf{(a)} Page~12 states that ``Table~2 presents the average protection decisions by
prior and signal type,'' but Table~2 actually groups observations by signal value
(white/black) and information-structure type (truthful, FP-only, FN-only, both
errors). The same mismatch appears in the caption of Table~A4.

\textbf{(b)} The authors could briefly justify why Tables~2, 3, and A4 pool across
priors. A single sentence would suffice.
\end{r2quote}

\response [Your response here.]

%----------------------------------------------------------------
\subsection{R2.3\quad Claim on Page~16 Regarding Introduction of the FN Rate}

\begin{r2quote}
\textbf{(a)} The paragraph below Table~3 claims ``a similar effect for the introduction
of the FN rate (row 1 vs.\ 2),'' but the mean updating error barely changes
(0.048 vs.\ 0.051).

\textbf{(b)} Moreover, introducing the FN rate changes the Bayesian posterior for the
white signal (from 0.000 to 0.119), so the row~1 vs.\ row~2 comparison does not hold
the posterior fixed � unlike row~1 vs.\ row~3.
\end{r2quote}

\response [Your response here.]

%----------------------------------------------------------------
\subsection{R2.4\quad Section~4 Closing Summary Understates the Evidence}

\begin{r2quote}
The paragraph before Section~5 concludes only that ``subjects understand the
experimental environment well enough.'' But Tables~2 and~3 also show that the FP (FN)
rate affects protection decisions and updating errors even when it should not influence
the Bayesian posterior � suggestive evidence for the FP/FN-confusion mechanism. A
sentence foreshadowing this in Section~4 would make the main mechanism visible
earlier.
\end{r2quote}

\response [Your response here.]

%----------------------------------------------------------------
\subsection{R2.5\quad Hypothesis~1 and Result~1 Overstate the t-Test Evidence}

\begin{r2quote}
Hypothesis~1 states WTPs ``are equal to their value for risk-neutral agents,'' and
Result~1 concludes ``no significant discrepancies.'' But the t-test only establishes
that $E[\Delta b_{is}]=0$; large but offsetting deviations could cancel. Suggested
revisions:
\begin{enumerate}[label=\roman*.]
  \item Restate more precisely: ``WTPs are not systematically above or below the
        risk-neutral benchmark \emph{on average}''; or
  \item Test the stronger claim by examining whether $|\Delta b_{is}|$ is significantly
        different from zero.
\end{enumerate}
\end{r2quote}

\response [Your response here.]

%----------------------------------------------------------------
\subsection{R2.6\quad (Optional) Reorganise Section~5.3 to Lead with FP/FN Confusion}

\begin{r2quote}
Given that FP/FN confusion is the primary mechanism, it might be preferable to present
it first and then evaluate and rule out alternative explanations.
\end{r2quote}

\response [Your response here.]

%----------------------------------------------------------------
\subsection{R2.7\quad Risk-Preference Discussion (pp.\ 23--24)}

\begin{r2quote}
The authors could note that the main effects of FP and FN costs remain significant,
rather than focusing solely on the insignificance of the interaction terms.
\end{r2quote}

\response [Your response here.]

%----------------------------------------------------------------
\subsection{R2.8\quad Probability Weighting Discussion (pp.\ 24--25)}

\begin{r2quote}
\textbf{(a)} In equation~(9) and the definition of $\Lambda$, the numerator should be
$P_{01}$ rather than $P_{00}$.

\textbf{(b)} The argument with rank-dependent utility is appreciated. The flow is clear
and the observation that $\pi + (1-\pi)P_{10}\leq 0.75$ is insightful. However, this
depends on the rotation point of the weighting function. If the rotation point lies far
to the left, that range might still fall in the steeper region, implying lower relative
sensitivity. Evidence for ``middle'' rotation points (e.g.\ Gonzalez and Wu, 1999)
would be helpful, though the flexibility of the weighting function may make this
difficult to rule out entirely.
\end{r2quote}

\response [Your response here.]

%----------------------------------------------------------------
\subsection{R2.9\quad Anticipation of Bias in Protection Decisions}

\begin{r2quote}
\textbf{(a)} Table~A10 should be Table~A8; the caption of Table~A8 should refer to WTP
\emph{deviations} rather than simply WTP.

\textbf{(b)} The two new variables ($\Delta\text{prot}$, $\Delta\text{loss}$) may not
adequately test the ``anticipation of bias'' mechanism. Negative coefficients suggest
that subjects who perform worse in the IP task are \emph{more} accurate in WTP, which
could indicate they anticipate their own mistakes.

\textbf{(c)} Given that the coefficients are statistically significant, concluding that
``biases in IP do not explain the observed pattern'' seems premature. A more appropriate
interpretation is that the coefficients on FP and FN do not \emph{change} after
introducing these variables.

\textbf{(d)} A more informative specification would include interaction terms between
mistake variables and FP/FN rates. If anticipation of bias is the primary channel,
mistakes should mediate the WTP--FP/FN relationship.

\textbf{(e)} Even if anticipation of bias is important, it does not undermine
FP/FN confusion; Tables~2 and~7 already show that IP decisions respond to FP/FN in
ways consistent with confusion. Anticipation of bias would reinforce, not replace, the
FP/FN-confusion mechanism.
\end{r2quote}

\response [Your response here.]

%----------------------------------------------------------------
\subsection{R2.10\quad Cross-Task Consistency Analysis}

\begin{r2quote}
The authors should consider incorporating the cross-task consistency analysis from the
response letter (Table~2, page~13) into the paper or at least the appendix.
\end{r2quote}

\response [Your response here.]

%----------------------------------------------------------------
\subsection{R2.11\quad (Optional) ``Dishonest Gremlin'' Framing}

\begin{r2quote}
Even without the ``dishonest'' framing, individuals may still conflate FP and FN errors.
A brief discussion of this potential limitation (perhaps at the end of the paper) would
be helpful.
\end{r2quote}

\response [Your response here.]

%----------------------------------------------------------------
\subsection{R2.12\quad Exposition}

\begin{r2quote}
\textbf{(a)} Page~11: ``this section also provides \ldots\ the theoretical predictions''
is imprecise � the predictions are those of the risk-neutral benchmark specifically.

\textbf{(b)} Relabel ``Posterior'' as ``Bayesian posterior'' in Tables~2, 3, A2, and A4.

\textbf{(c)} In Tables~5 and~A8, a horizontal rule separating the main regression panel
from the risk-averse/risk-loving subgroup tests would improve readability.

\textbf{(d)} (Optional.) Adding a short closing sentence to each appendix table caption
pointing to the key coefficient or column would reduce cross-referencing costs.

\textbf{(e)} Page~21: ``the systematic WTP biases reported above'' is vague � please
specify which biases are meant.
\end{r2quote}

\response [Your response here.]

%----------------------------------------------------------------
\subsection{R2.13\quad Typos}

\begin{r2quote}
\begin{itemize}[noitemsep]
  \item Equation~(2), page~5: extra ``$=$'' sign.
  \item Page~7: ``Equations 6-s 7'' $\to$ ``Equations 6 and 7.''
  \item Footnotes~10 and~11: missing full stop at end.
  \item Footnote~13, page~21: ``identical'' should likely be ``different.''
  \item Last paragraph, page~26: ``Table~A10'' $\to$ ``Table~A8.''
\end{itemize}
\end{r2quote}

\response All five typos have been corrected. We thank the referee for the careful
reading.

%----------------------------------------------------------------
\end{document}